% =============================================================================
% Claude for Financial Services — minimal anchor deck (hybrid delivery)
% 4 slides, designed to be brought up at specific beats in a discussion-format
% small-room talk (≤12 people). The full deck lives in
% claude-financial-services.tex and is the backup / handout / reference.
%
% Compile:
%   latexmk -pdf -xelatex claude-financial-services-minimal.tex
% =============================================================================
\documentclass[11pt,aspectratio=169]{beamer}

% --- Theme -------------------------------------------------------------------
\usetheme{metropolis}
\metroset{
  progressbar=none,
  numbering=none,
  block=fill,
  sectionpage=none,
  subsectionpage=none
}

% --- Encoding & fonts --------------------------------------------------------
\usepackage[english]{babel}
\usepackage{microtype}
\usepackage{fontspec}

% --- Colour palette (same as full deck) --------------------------------------
\definecolor{accentpurple}{HTML}{7C3AED}
\definecolor{darkgray}{HTML}{1F2937}
\definecolor{lightgray}{HTML}{F3F4F6}
\definecolor{successgreen}{HTML}{059669}
\definecolor{warningred}{HTML}{DC2626}
\setbeamercolor{frametitle}{bg=darkgray}
\setbeamercolor{title separator}{fg=accentpurple}

% --- Tables / graphics -------------------------------------------------------
\usepackage{booktabs}
\usepackage{array}
\usepackage{xcolor}

% hyperref is auto-loaded by beamer
\hypersetup{hidelinks}

% --- Custom commands ---------------------------------------------------------
\newcommand{\good}{\textcolor{successgreen}{\checkmark}}
\newcommand{\bad}{\textcolor{warningred}{\textbf{\texttimes}}}
\newcommand{\hl}[1]{\textbf{\textcolor{accentpurple}{#1}}}

% =============================================================================
\begin{document}

% --- Slide 1: The headline ---------------------------------------------------
\begin{frame}{Claude for Financial Services --- the headline}
  On \hl{5 May 2026}, Anthropic launched Claude for Financial Services:
  \begin{itemize}
    \item \textbf{10 ready-to-run agent templates} for high-volume FS work
    \item \textbf{Microsoft 365 integration} (Excel, PowerPoint, Word, Outlook)
    \item \textbf{MCP-based connectors} to $\sim$11 financial data providers
    \item \textbf{Enterprise controls}: SSO, SCIM, audit logs, ISO/IEC 42001:2023
  \end{itemize}
  \vfill
  Underlying model: \hl{Claude Opus 4.7} --- leads Vals AI Finance Agent benchmark at \hl{64.37\%}.
\end{frame}

% --- Slide 2: The 10 agent templates -----------------------------------------
\begin{frame}{The 10 agent templates}
  \centering\footnotesize
  \begin{tabular}{lp{0.6\textwidth}}
    \toprule
    \multicolumn{2}{l}{\textit{Research \& Client Coverage}} \\
    \midrule
    Pitch Builder       & Target lists $\rightarrow$ comps $\rightarrow$ pitchbook + cover note \\
    Meeting Preparer    & Client / counterparty briefs ahead of calls \\
    Earnings Reviewer   & Transcripts + filings $\rightarrow$ model updates + flagged changes \\
    Model Builder       & Builds models from filings, feeds, analyst inputs \\
    Market Researcher   & Sector / issuer tracking; news + filings + broker synthesis \\
    \midrule
    \multicolumn{2}{l}{\textit{Finance \& Operations}} \\
    \midrule
    Valuation Reviewer  & Checks valuations vs comparables and firm standards \\
    GL Reconciler       & Reconciles GL accounts; NAV calculations \\
    Month-End Closer    & Close checklist, journal entries, close reports \\
    Statement Auditor   & Statement consistency and audit-readiness review \\
    KYC Screener        & Assembles entity files, packages escalations \\
    \bottomrule
  \end{tabular}
\end{frame}

% --- Slide 3: What we have / don't have --------------------------------------
\begin{frame}{What we have --- and don't}
  \begin{columns}[T,onlytextwidth]
    \begin{column}{0.46\textwidth}
      \textbf{Have} \medskip
      \begin{itemize}
        \item \good\ Microsoft 365 Copilot
        \item \good\ GitHub Copilot
      \end{itemize}
      \medskip
      \footnotesize\textit{The full list of sanctioned office AI tooling.}
    \end{column}
    \begin{column}{0.5\textwidth}
      \textbf{Don't have} \medskip
      \begin{itemize}
        \footnotesize
        \item \bad\ Claude.ai (consumer)
        \item \bad\ Claude Team / Enterprise
        \item \bad\ Claude for M365 plugin
        \item \bad\ Claude Cowork
        \item \bad\ Claude Code
        \item \bad\ Anthropic API (org-level)
      \end{itemize}
    \end{column}
  \end{columns}
  \vfill
  \hl{No path} from M365 Copilot or GitHub Copilot to the Claude FS agents. The two product families are separate.
\end{frame}

% --- Slide 4: What to tell colleagues ----------------------------------------
\begin{frame}{What to tell colleagues asking about Claude}
  \begin{exampleblock}{Clean two-sentence version}
    \textit{Yes --- Anthropic launched Claude for Financial Services in May. It's a strong product for our type of work, but we don't have a sanctioned channel to use it: our office AI tooling today is Microsoft 365 Copilot and GitHub Copilot, and Claude requires a separate subscription we haven't acquired.}
  \end{exampleblock}

  \vfill
  \begin{block}{If they push: ``can I just use my personal account?''}
    \textit{Not on office data, no --- that would be a data-handling issue regardless of how useful the output is. Personal Claude is fine for personal exploration; office work goes through office-sanctioned tools.}
  \end{block}
\end{frame}

\end{document}
